\documentclass[12pt]{article}


\usepackage[upright]{fourier}% l'option permet d'avoir les majuscules droites dans les formules mathématiques
\usepackage[utf8x]{inputenc}
\usepackage[french]{babel}
\AddThinSpaceBeforeFootnotes % à insérer si on utilise \usepackage[french]{babel}
\FrenchFootnotes % à insérer si on utilise \usepackage[french]{babel}
%s'utilise ainsi : Les Représentants\footnote{ceci est une note} du Peuple Français
\DecimalMathComma %supprime l'espace après la virgule dans un nombre
\usepackage{enumerate} % pour pouvoir changer la numérotation des listes
\usepackage[scaled=0.875]{helvet}
\renewcommand{\ttdefault}{lmtt}
\usepackage{mathtools,amsmath,amssymb,makeidx,bm} %bm : pour mettre en gras dans une formule de maths

%packages permettant d'augmenter le nombre de registres de dimension et donc d'éviter les erreurs de compilation dûs aux packages tikz, pstricks and co
\usepackage{etex}



%redéfinition de fractions, limites, sommes, intégrales, coefficients binomiaux en displaystyle, limites de suites
\let\binomOld\binom
\renewcommand{\binom}{\displaystyle\binomOld}
\let\limOld\lim
\renewcommand{\lim}{\displaystyle\limOld}
\newcommand{\limn}{\lim_{n\to +\infty}}
\newcommand{\limm}{\lim_{x\to -\infty}}
\newcommand{\limp}{\lim_{x\to +\infty}}
\newcommand{\limz}{\lim_{x\to 0}}
\newcommand{\limzm}{\lim_{\substack{x \to 0\\ x < 0}}}
\newcommand{\limzp}{\lim_{\substack{x \to 0\\ x > 0}}}
\let\sumOld\sum
\renewcommand{\sum}{\displaystyle\sumOld}
\let\intOld\int
\renewcommand{\int}{\displaystyle\intOld}

%Calligraphie spéciale
\usepackage{mathrsfs}
\newcommand{\scr}{\mathscr}

%ssi
\newcommand{\lr}{\leftrightarrow}
\newcommand{\llr}{\longleftrightarrow}
\newcommand{\Lr}{\Leftrightarrow}
\newcommand{\LLr}{\Longleftrightarrow}

%Nombres complexes
\let\Reold\Re
\renewcommand{\Re}{~\text{Re}~}
\let\Imold\Im
\renewcommand{\Im}{~\text{Im}~}
\newcommand{\ii}{\,\mathrm{i}}

%style des pages
\usepackage{fancybox,cancel}
\usepackage{fancyhdr}
\usepackage{lastpage}
%redéfinition du style plain
\fancypagestyle{plain}{%
\fancyhf{} %vide l'en-tête et le pied de page
%\fancyfoot[C]{\bfseries \thepage / \pageref{LastPage}} %numéro de la page en gras et centré
\fancyfoot[C]{\thepage / \pageref{LastPage}} %numéro de la page centré
\renewcommand{\headrulewidth}{0.2pt}
\renewcommand{\footrulewidth}{0.2pt}}

%redéfinition des marges par défaut
\setlength\paperheight{297mm}
\setlength\paperwidth{210mm}
\setlength{\evensidemargin}{0cm}
\setlength{\oddsidemargin}{-0.5cm}
\setlength{\topmargin}{-2cm}
\setlength{\headsep}{0cm}
\setlength{\headheight}{0.7cm}
\setlength{\textheight}{25.2cm}
\setlength{\textwidth}{17cm}

%TABLEAU
%diminuer la taille des caractères dans un tableau
%Pour aérer un tableau, il est possible de redéfinir l'espacement entre les lignes d'un tableaux et l'espacement entre les colonnes, par exemple :
\usepackage{array}
%\setlength{\tabcolsep}{1cm}
\renewcommand{\arraystretch}{1.5}%augmente la hauteur des lignes des tableaux
%colonnes centrées verticalement et horizontalement permettant d'écrire des paragraphes de largeur fixée du type M{3cm}
\newcolumntype{M}[1]{>{\centering\arraybackslash}m{#1}}%\arraybackslash permet de continuer à utiliser \\ pour le changement de ligne
%couleurs cellules, colonnes, lignes
\usepackage{color,colortbl}
\usepackage[table]{xcolor}
\usepackage{longtable}%permet d'obtenir des tableaux sur plusieurs pages
%\usepackage{arydshln}% pour pouvoir ajouter des lignes horizontales en pointillés avec \hdashline au lieu de \hline
%pose des problèmes d'encadré avec alterqcm


%ALGORITHME avec Algobox
\usepackage{ucs}
%\usepackage{xcolor}
\usepackage{framed}
\definecolor{fond}{gray}{0.95}
\newenvironment{cadrecode}{%
  \def\FrameCommand{{\color[HTML]{888888}\vrule width 3pt}\colorbox{fond}}%
  \MakeFramed {\advance\hsize-\width \FrameRestore}}%
{\endMakeFramed}
\usepackage{alltt}

%arbres
\usepackage{pstricks,pst-plot,pst-text,pst-tree,pstricks-add,pst-eucl,pst-all}

%tabvar %packages à ajouter aux précédents
\usepackage{pst-eps,pst-fill,pst-node,pst-math}
%\input tabvar %ajouter le fichier tabvar.tex dans le même dossier que le fichier actuel
%à utiliser uniquement lors d'un tableau de signes et de variations

%INTERVALLES
\newcommand{\intervalleff}[2]{\left[#1\,;#2\right]}
\newcommand{\intervallefo}[2]{\left[#1\,;#2\right[}
\newcommand{\intervalleof}[2]{\left]#1\,;#2\right]}
\newcommand{\intervalleoo}[2]{\left]#1\,;#2\right[}


%INTERLIGNES
\usepackage{setspace}
%s'utilise avec \begin{spacing}{''facteur''}
%   […]
%\end{spacing}

%IMAGES
\usepackage{graphicx} %inclure des graphiques
\usepackage{caption}
\usepackage{subfig}

\usepackage{tabularx}
\usepackage{soul} % Pour souligner : \ul
\usepackage{ulem} % Pour souligner double : \uuline
                      % Pour souligner ondulé : \uwave
                      % Pour barrer horizontal : \sout
                      % Pour barrer diagonal : \xout
\usepackage{pifont}
\usepackage{slashbox}
\usepackage{textcomp}
\usepackage{pst-plot}

\usepackage[np]{numprint}


%ENCADRES
\usepackage{pstricks}
\usepackage{pst-grad}
\usepackage{xkeyval}
\usepackage{pst-coil}
\usepackage{ifthen}
\usepackage{ifpdf}
\usepackage{pst-blur}
\usepackage{bclogo}
%doc ici :
%http://www.tug.org/texlive/Contents/live/texmf-dist/doc/latex/bclogo/bclogo-doc.pdf

%figures tikz
\usepackage{pgf,tikz}
\usetikzlibrary{arrows}


%raccourcis perso
\newcommand\pfr[1]{\psframebox[linecolor=red]{#1}}
\newcommand\coef[1][]{c{\oe}fficient#1\xspace}
\newcommand\abs[1]{\ensuremath{\left\vert #1 \right\vert}}

%liens hypertexte
\usepackage[colorlinks=true,linkcolor=black,filecolor=blue,urlcolor=blue]{hyperref}
%\href{mailto:mail@exemple.com}{Mon mail} % Lien email
%\href{http://exemple.com}{Mon site web}  % Lien web
%\href{fichier.pdf}{Mon fichier}          % Lien vers un fichier

\newcommand{\euro}{\texteuro{}~}

%simplification notation norme \norme{}
\newcommand{\norme}[1]{\left\Vert #1\right\Vert}

%Ensembles de nombres
\usepackage{dsfont}
\newcommand{\C}{\mathds C}
\newcommand{\R}{\mathds R}
\newcommand{\Q}{\mathds Q}
\newcommand{\D}{\mathds D}
\newcommand{\Z}{\mathds Z}
\newcommand{\N}{\mathds N}
\newcommand{\e}{\mathrm {e}}
\newcommand{\dd}{\,\mathrm{d}}



%simplification de la notation de vecteur \vect{}
\newcommand{\vect}[1]{\mathchoice%
{\overrightarrow{\displaystyle\mathstrut#1\,\,}}%
{\overrightarrow{\textstyle\mathstrut#1\,\,}}%
{\overrightarrow{\scriptstyle\mathstrut#1\,\,}}%
{\overrightarrow{\scriptscriptstyle\mathstrut#1\,\,}}}


\renewcommand{\theenumi}{\textbf{\arabic{enumi}}}
\renewcommand{\labelenumi}{\textbf{\theenumi.}}
\renewcommand{\theenumii}{\textbf{\alph{enumii}}}
\renewcommand{\labelenumii}{\textbf{\theenumii.}}

%Arc de cercles
\newlength{\longarc}
\newcommand{\arc}[1]{\settowidth{%
\longarc}{$#1$}
\addtolength{\longarc}{-0.5em}%
\unitlength \longarc \ensuremath{%
\stackrel{\begin{picture}(1,0.2)
\qbezier(0,0)(0.5,0.2)(1,0)
\end{picture}}{#1}}}
%l'arc AB s'obtient par \arc{AB} en mode mathématiques ou non
%autre possibilité avec deux lettres : \overset{\displaystyle\frown}{AB}


%Repères
\def\Oij{$\left(\text{O}~;~\vect{\imath},~\vect{\jmath}\right)$}
\def\Oijk{$\left(\text{O}~;~\vect{\imath},~ \vect{\jmath},~ \vect{k}\right)$}
\def\Ouv{$\left(\text{O}~;~\vect{u},~\vect{v}\right)$}


%Symboles Casio dans les programmes
\usepackage{pifont} % ça c'est pour le retour chariot
\newcommand{\RetourChariot}{\Pisymbol{psy}{191}}

%Pointillés sur toute la ligne
\usepackage{multido}

\newcommand{\Pointilles}[1][1]{%
\multido{}{#1}{\makebox[\textwidth]{\dotfill}\\[\parskip]
}}
%commandes : \Pointilles ou \Pointilles[4] pour 4 lignes


%textes à trous
\newlength\lgtrou
\newcommand*\trou[1]{%
\settowidth\lgtrou{#1}%
\hspace*{2\lgtrou}}
\setlength\baselineskip{1.2\baselineskip}
%Commande à utiliser : \trou{texte qui sera remplacé par un espace vide}



%%%% Papier millim\'er\'ex*y en cm
\newcommand{\PapierMill}[2]{%
	\begin{pspicture*}(0,0)(#1,#2)
	
	\psgrid[gridwidth=0.3pt,
		gridlabels=0,
		subgriddiv=10,
		subgridcolor=black,
		subgridwidth=0.05pt](0,0)(21,29.7) % il faut \'ere sr
				% que \'e d\'easse, sinon il y a des soucis avec
				% des lignes qui ne vont pas jusqu'au bout
				% \'ecause des arrondis entier de multido
				% pspicture* coupe aux bonnes dimensions
	
	% Grille de 5cm en 5cm
	\psgrid[xunit=5,yunit=5,gridwidth=0.5pt,
		gridlabels=0,
		subgriddiv=0](0,0)(#1,#2)%
		
	\psline[linewidth=1.5pt](0,0)(0,#2)(#1,#2)(#1,0)(0,0)	
	\end{pspicture*}
	}



% Ecrire sur plusieurs colonnes
\usepackage{multicol}
\setlength{\columnseprule}{0.5pt}
%exemple
%\begin{multicols}{3}[Titre sur une seule colonne.]
%   3~colonnes équilibrées, 3~colonnes équilibrées, 3~colonnes équilibrées, 3~colonnes équilibrées
%\end{multicols}
%\begin{multicols}{2}[\section{Titre numéroté.}]
%   blabla sur deux colonnes, c'est plus sérieux. C'est le style qui est généralement utilisé pour écrire des articles.
%saut de colonne forcé : 
%\columnbreak
%djhskjdhjsq
%sdkksqjhd
%\end{multicols}
%Pour ajouter un titre numéroté qui apparaisse sur toute la largeur de la page, il faut utiliser l'option [\section{Titre.}] juste après \begin{multicols}{nb-col}.
%Remarques :
%Pour qu'une ligne de séparation apparaisse entre les colonnes, il faut utiliser : \setlength{\columnseprule}{1pt}.

%Pour redéfinir la largeur de l'espace inter-colonnes, il faut utiliser \setlength{\columnsep}{30pt}.

%Pour remonter le texte, dans chaque colonne vers le haut : \raggedcolumns qui se tape :\begin{multicols}{2}\raggedcolumns...\columnbreak...\columnbreak\end{multicols}

%Pour supprimer les traits verticaux : \setlength{\columnseprule}{0pt} avant \begin{multicols}{3}...\end{multicols}


%QRcode, codebarre
\usepackage{pst-barcode}
%\begin{pspicture}(2,2)
%	\psbarcode{http://www.latex-howto.be}{eclevel=M}{qrcode}
%\end{pspicture}

%Texte en filigrane
\usepackage{watermark}
%On utilise ensuite les commandes \watermark, \leftwatermark, \rightwatermark ou \thiswatermark qui permettent de définir un filigrane sur toutes les pages, les pages paires, les pages impaires ou juste une page
%Exemple : \thiswatermark {
%\begin{minipage}{0.95\linewidth}
%\vspace{25cm}
%\begin{center}
%\rotatebox{55}{\scalebox{8}{\color[gray]{0.7}\LaTeX}}
%\end{center}
%\end{minipage}
%}


%Numérotation des sections
\renewcommand{\thesection}{\Roman{section}\xspace}
\renewcommand{\thesubsection}{\thesection .\Alph{subsection}\xspace} %lettres majuscules
\renewcommand{\thesubsubsection}{\thesubsection .\alph{subsubsection}\xspace} %lettres minuscules


% Mise en forme des algorithmes
\usepackage[french,boxed,titlenumbered,lined,longend]{algorithm2e}
  \SetKwIF {Si}{SinonSi}{Sinon}{si}{alors}{sinon\_si}{alors}{fin~si}
 \SetKwFor{Tq}{tant\_que~}{~faire~}{fin~tant\_que}
 \SetKwFor{PourCh}{pour\_chaque }{ faire }{fin pour\_chaque}
 \SetKwInput{Sortie}{Sortie}
  \SetKwInput{Entree}{Entrée}
\newcommand{\Algocmd}[1]{\textsf{\textsc{\textbf{#1}}}}\SetKwSty{Algocmd}
  \newcommand{\AlgCommentaire}[1]{\textsl{\small  #1}} 

%symbole parallèle avec \sslash
\usepackage{stmaryrd}


%%%%%%%%%%%%%%%%%%%%%%%%%%%%%%%%%%%%%%%%%%%%%%%%%%%%%%%%%%%%%%%%%%%%%%%%%%%%%%%
%%%%%%%%%%%%%%%%%%%%%%%%%%%%%%%%%%%%%%%%%%%%%%%%%%%%%%%%%%%%%%%%%%%%%%%%%%%%%%%
%Encadrés pour Propriétés, Théorème, Définitions, exemples, exercices
\newcounter{propr}
\newenvironment{propr}{\refstepcounter{propr}\begin{bclogo}[couleur = gray!30 , arrondi = 0.1 ,logo = \bclampe , barre = snake , tailleOndu = 1.5]{Propriété \thepropr}}{\end{bclogo}}
\newcounter{theo}
\newenvironment{theo}{\refstepcounter{theo}\begin{bclogo}[couleur = gray!50 , arrondi = 0.1 ,logo = \bclampe , barre = snake , tailleOndu = 1.5]{Théorème \thetheo}}{\end{bclogo}}
\newcounter{defi}
\newenvironment{defi}{\refstepcounter{defi}\begin{bclogo}[couleur = blue!30 , arrondi = 0.1 ,logo=\bclampe , barre = snake , tailleOndu = 1.5]{Définition \thedefi}}{\end{bclogo}}

\newcounter{exemple}\newcommand{\exemple}{\refstepcounter{exemple}\textbf{Exemple \theexemple \ :}}
\newcounter{exercice}\newcommand{\exercice}{\refstepcounter{exercice}\textbf{\large{Exercice \theexercice \ :}}}
\newcounter{remarque}\newcommand{\remarque}{\refstepcounter{remarque}\textbf{Remarque \theremarque \ :}}


%%%%%%%%%%%%%%%%%%%%%%%%%%%%%%%%%%%%%%%%%%%%%%%%%%%%%%%%%%%%%%%%%%%%%%%%%%%%%%%

%QCM
\usepackage{alterqcm}					%%Permet de créer des QCM
%\begin{alterqcm}
%\AQquestion{Question}{{Proposition 1},{Proposition 2},{Proposition 3}}
%\end{alterqcm}


%%%%%%%%%%%%%%%%%%%%%%%%%%%%%%%%%%%%%%%%%%%%%%%%%%%%%%%%%%%%%%%%%%%%%%%%%%%%%%%






%%%%%%%%%%%%%%%%%%%%%%%%%%%%%%%%%%%%%%%%%%%%%%%%%%%%%%%%%%%%%%%%%%%%%%%%%%%%%%%

%entête classique

\newcommand{\tete}{
	{\noindent
      \textbf{BAT 2}     	            %% Auteur
      \hfill \textbf{Année \np{2013}-\np{2014}}   \\  	            	%% Année
	}
	%{\noindent
  %    \url{tremulot.maths@gmail.com}                 	%% Module
  %    \hfill \textbf{\today}   \\  % Numero de la feuille
	%}
	\hrule
	\begin{center} 
	\Large \textsc{\textbf{Chapitre 1 : Nombres complexes}}
	\end{center}
	\hrule
	\noindent 
\medskip}
%%%%%%%%%%%%%%%%%%%%%%%%%%%%%%%%%%%%%%%%%%%%%%%%%%%%%%%%%%%%%%%%%%%%%%%%%%%%%%%

%entête contrôle

\newcommand{\controle}{
\begin{center}
\begin{tabularx}{\textwidth}{|c|X|c|}
%{@{\extracolsep{\fill}} |l c r|}
		\hline
		BAT 2 & \centering Contrôle de mathématiques n° 1 (durée : 55 minutes) & Année \np{2013}/\np{2014} \\
		\hline
		\multicolumn{3}{|c|}{Mercredi 24 avril} \\
		\hline
		\multicolumn{3}{|l|}{NOM, Prénom : \dotfill} \\
		\hline
\end{tabularx}

\vspace{1cm}

%\textbf{Sujet B}
%\end{center}
%\begin{center}
\textit{La qualité et la précision de la rédaction seront prises en compte dans l'appréciation des copies.\\ Faites des phrases simples et précises.\\
Les calculatrices électroniques de poche sont autorisées, conformément à la règlementation en vigueur.}
\end{center}

\noindent
%\vspace{1cm}
}

%%%%%%%%%%%%%%%%%%%%%%%%%%%%%%%%%%%%%%%%%%%%%%%%%%%%%%%%%%%%%%%%%%%%%%%%%%%%%%%

%entête DM

\newcommand{\dm}{
\begin{minipage}{0.9\linewidth}
\begin{center}
\begin{tabularx}{\textwidth}{|c|X|c|}
%{@{\extracolsep{\fill}} |l c r|}
		\hline
		BAT 2 & \centering Devoir maison de mathématiques n° 1 & Année \np{2013}/\np{2014} \\
		\hline
		\multicolumn{3}{|c|}{\`A rendre, sur feuille, au plus tard le 15/05} \\
		\hline
		\multicolumn{3}{|l|}{NOM, Prénom : \dotfill} \\
		\hline
\end{tabularx}
\end{center}
\end{minipage}
\hfill
\begin{minipage}{0.09\linewidth}
\begin{pspicture}(2,2)
	\psbarcode{http://www.tremulotmaths.fr}{eclevel=L}{qrcode}
\end{pspicture}
\end{minipage}

\begin{center}
\textbf{Vous pouvez poser vos questions dans les commentaires de l'\href{http://www.tremulotmaths.fr/spip.php?article308}{article sur le site}.}
\end{center}
\noindent
}
%%%%%%%%%%%%%%%%%%%%%%%%%%%%%%%%%%%%%%%%%%%%%%%%%%%%%%%%%%%%%%%%%%%%%%%%%%%%%%%%%%%%%%
%entête bac blanc
\newcommand{\bacblanc}{
\thispagestyle{empty}
\thiswatermark {
\begin{minipage}{0.95\linewidth}
\vspace{25cm}
\begin{center}
\rotatebox{55}{\scalebox{8}{\color[gray]{0.7}OBLIGATOIRE}}
\end{center}
\end{minipage}
}
Lycée Jean-Pierre Timbaud \hfill Février \np{2014}

.\dotfill
\vspace{2cm}
\begin{center}
\textbf{\LARGE{SESSION \np{2014}}}\\
\vspace{1cm}
\textbf{\LARGE{BACCALAUR\'EAT G\'EN\'ERAL BLANC}}\\
\vspace{1cm}
\textbf{\LARGE{Série SCIENTIFIQUE}}\\
\vspace{1cm}
\textbf{\LARGE{MATH\'EMATIQUES}}

\end{center}

\vspace{1cm}

Durée : 4 heures 
\hfill
Coefficient 7

\vspace{1cm}

\dotfill
\begin{center}
L'usage des calculatrices est autorisé pour cette épreuve.\\
\end{center}
\dotfill

\vspace{1cm}


\begin{center}
\begin{bf}

Le candidat doit traiter les quatre exercices.\\
Le candidat est invité à faire figurer sur la copie toute trace de recherche, même incomplète ou non fructueuse, qu'il aura développée.\\
\vspace{12pt}
Il est rappelé que la qualité de la rédaction, la clarté et la précision des raisonnements entreront pour une part importante dans l'appréciation des copies.\\
\vspace{12pt}
Dès que le sujet vous est remis, assurez-vous qu'il est complet, que toutes les pages sont imprimées.\\
Le sujet comporte 5 pages.\\

\end{bf}
\end{center}
\newpage
}
%%%%%%%%%%%%%%%%%%%%%%%%%%%%%%%%%%%%%%%%%%%%%%%%%%%%%%%%%%%%%%%%%%%%%%%%%%%%%%%

%%%%%%%%%%%%%%%%%%%%%%%
%% DEBUT DU DOCUMENT %%
%%%%%%%%%%%%%%%%%%%%%%%

\begin{document}

\setlength\parindent{0mm}
%\tete 		%entête classique
%\controle		%entête contrôle
\dm 		%entête DM
%\enlargethispage{2cm}
%\bacblanc 		%entête bac blanc

\renewcommand \footrulewidth{0.2pt}%
\renewcommand \headrulewidth{0pt}%
\pagestyle{fancy}
\fancyhf{}
\cfoot{\thepage / \pageref{LastPage}}
\lfoot{BAT 2 - Nombres complexes}
\rfoot{Année \np{2013}-\np{2014}}

%\thispagestyle{empty} %annuler le style pour la page en cours, soit la première page, ici


%\title{BACCALAUR\'EAT BLANC}
%\date{}%supprime la date si on n'entre rien entre les accolades

%\maketitle %impose le style plain à la page
%\hrule

%%%%%%%%%%%%%%%%%%%%%%%Page 1%%%%%%%%%%%%%%%%%%%%%%%%%%%%%%%%%%%%%%%%
\begin{spacing}{1.2}

%Modèle propriété :
%\fbox{%
%   \begin{minipage}{\textwidth}
%   \textbf{\ul{Définition :}} \\
%   L'ensemble des matrices d'ordre $(n,p)$ à coefficients réels se note $\mathcal{M}_{n,p}(\R)$.
%   \end{minipage}%
%}


%Autre possibilité avec les packages de la partie ENCADRES
%\begin{bclogo}[couleur = gray!30 , arrondi = 0.1 ,logo = \bclampe , barre = snake , tailleOndu = 1.5]{Théorème}
%Soit $M$ la matrice de transition d'un graphe probabiliste.
%Soit $P_0$ la matrice ligne décrivant l'état initial.
%La matrice ligne $P_n$ décrivant l'état probabiliste à l'étape $n$ vérifie :
%$\bullet P_{n+1}=P_n\times M$
%$\bullet P_n=P_0\times M^n$
%\end{bclogo}
%%%%%%%%%%%%%%%%%%%%%%%%%%%%%%%%%%%%%%%%%%%%%%%%%%%%%%%%%%%%
\renewcommand{\contentsname}{Sommaire}\tableofcontents \newpage
%%%%%%%%%%%%%%%%%%%%%%%%%%%%%%%%%%%%%%%%%%%%%%%%%%%%%%%%%%%%
\thiswatermark {
\begin{minipage}{0.95\linewidth}
\vspace{25cm}
\begin{center}
\rotatebox{55}{\scalebox{8}{\color[gray]{0.7}\LaTeX 2}}
\end{center}
\end{minipage}
}

\section{Introduction}

\subsection{Identification de la situation}

\begin{center}
\begin{pspicture}(2,2)
	\psbarcode{http://www.tremulotmaths.fr}{eclevel=M}{qrcode}
\end{pspicture}

\end{center}
On se situe, dans ce chapitre, dans deux domaines des statistiques, qui sont ceux de \og l'échantillonnage\fg~ et de \og l'estimation\fg~.

Une ligne de pointillés :

\Pointilles

Trois lignes de pointillés :

\Pointilles[3]

\begin{minipage}{0.45\linewidth}

Une ligne de pointillés dans minipage :
\Pointilles

\end{minipage}
\hfill
\begin{minipage}{0.45\linewidth}

Trois lignes de pointillés dans minipage :
\Pointilles[3]

\end{minipage}

Apprenons à reconnaître les différents contextes d'applications.

\medskip


\textbf{\textit{Exemples :}}

On considère deux urnes $U_1$ et $U_2$ contenant chacune un très grand nombre de boules, rouges ou bleues.

\begin{enumerate}
\item 
Dans l'urne $U_1$, \textbf{on connaît la proportion $p$ de boules rouges}.

\begin{minipage}{0.45\linewidth}
On procède à des tirages avec remise de $n$ boules, et on 

observe la fréquence d'apparition d'une boule rouge.

Cette fréquence observée appartient \og en général\fg~ à un \og intervalle de fluctuation \fg~ de centre $p$, dont la longueur diminue lorsque $n$ augmente.

\end{minipage}
\hfill
\begin{minipage}{0.45\linewidth}
\includegraphics[width=5cm]{fig.eps}

\end{minipage}


Cet intervalle est un \og intervalle de fluctuation \fg~.

On est ici dans le domaine de \textbf{l'échantillonnage} et de l'\textbf{intervalle de fluctuation}.
\[\intervalleff{\dfrac{5}{2}}{\sqrt{\dfrac{1}{3}}}\]
\[\intervallefo{\dfrac{5}{2}}{\sqrt{\dfrac{1}{3}}}\]
\[\intervalleof{\dfrac{5}{2}}{\sqrt{\dfrac{1}{3}}}\]
\[\intervalleoo{\dfrac{5}{2}}{\sqrt{\dfrac{1}{3}}}\]

\begin{center}
\definecolor{cqcqcq}{rgb}{0.75,0.75,0.75}
\begin{tikzpicture}[line cap=round,line join=round,>=triangle 45,x=1.0cm,y=1.0cm]
\draw [color=cqcqcq,dash pattern=on 2pt off 2pt, xstep=1.0cm,ystep=1.0cm] (-0.46,-3.28) grid (6.88,3.36);
\draw[->,color=black] (-0.46,0) -- (6.88,0);
\foreach \x in {,1,2,3,4,5,6}
\draw[shift={(\x,0)},color=black] (0pt,2pt) -- (0pt,-2pt) node[below] {\footnotesize $\x$};
\draw[->,color=black] (0,-3.28) -- (0,3.36);
\foreach \y in {-3,-2,-1,1,2,3}
\draw[shift={(0,\y)},color=black] (2pt,0pt) -- (-2pt,0pt) node[left] {\footnotesize $\y$};
\draw[color=black] (0pt,-10pt) node[right] {\footnotesize $0$};
\clip(-0.46,-3.28) rectangle (6.88,3.36);
\draw[smooth,samples=100,domain=0.001:6.88] plot(\x,{ln((\x))});
\end{tikzpicture}
\end{center}

\item Dans l'urne $U_2$, \textbf{on ignore la proportion de boules rouges}.

En procédant à des tirages avec rmise de $n$ boules, on va essayer d'estimer la proportion $p$ de boules rouges dans l'urne.

Cette estimation se fait au moyen d'un \og intervalle de confiance \fg~.

Cet intervalle dépend d'un coefficient, le \og niveau de confiance\fg~, que l'on attribue à l'estimation.

On est ici dans le domaine de l'\textbf{estimation}, et de l'\textbf{intervalle de confiance}.

\begin{center}
\definecolor{cqcqcq}{rgb}{0.75,0.75,0.75}
\begin{tikzpicture}[line cap=round,line join=round,>=triangle 45,x=1.0cm,y=1.0cm]
\draw [color=cqcqcq,dash pattern=on 4pt off 4pt, xstep=2.0cm,ystep=2.0cm] (-1.95,-1.49) grid (8.29,6.99);
\draw[->,color=black] (-1.95,0) -- (8.29,0);
\foreach \x in {,2,4,6,8}
\draw[shift={(\x,0)},color=black] (0pt,2pt) -- (0pt,-2pt) node[below] {\footnotesize $\x$};
\draw[->,color=black] (0,-1.49) -- (0,6.99);
\foreach \y in {,2,4,6}
\draw[shift={(0,\y)},color=black] (2pt,0pt) -- (-2pt,0pt) node[left] {\footnotesize $\y$};
\draw[color=black] (0pt,-10pt) node[right] {\footnotesize $0$};
\clip(-1.95,-1.49) rectangle (8.29,6.99);
\draw[smooth,samples=100,domain=0.01:8.29] plot(\x,{ln((\x))});
\draw[smooth,samples=100,domain=0:1.95] plot(\x,{2.718^(\x)});
\draw[smooth,samples=100,domain=-1.95:0] plot(\x,{2.718^(\x)});
\draw[smooth,samples=100,domain=-1.95:8.29] plot(\x,{(\x)});
\draw (5.07,3.19) node[anchor=north west] {$y=\ln (x)$};
\draw (4.01,6.02) node[anchor=north west] {$y=x$};
\draw (1.14,5.28) node[anchor=north west] {$y=\text{e}^x$};
\end{tikzpicture}

\psset{xunit=2.0cm,yunit=2.0cm}
\begin{pspicture*}(-0.98,-0.74)(4.14,3.49)
\psgrid[subgriddiv=0,gridlabels=0,gridcolor=lightgray](0,0)(-0.98,-0.74)(4.14,3.49)
\psset{xunit=1.0cm,yunit=1.0cm,algebraic=true,dotstyle=o,dotsize=3pt 0,linewidth=0.8pt,arrowsize=3pt 2,arrowinset=0.25}
\psaxes[labelFontSize=\scriptstyle,xAxis=true,yAxis=true,Dx=2,Dy=2,ticksize=-2pt 0,subticks=2]{->}(0,0)(-1.95,-1.49)(8.29,6.99)
\psplot[plotpoints=200]{1.1299160078328312E-8}{8.289582400000004}{ln(x)}
\psplot[plotpoints=200]{-1.953261199999997}{8.289582400000004}{EXP(x)}
\psplot[plotpoints=200]{-1.953261199999997}{8.289582400000004}{x}
\rput[tl](5.07,3.19){$y=\ln (x)$}
\rput[tl](4.01,6.02){$y=x$}
\rput[tl](1.14,5.28){$y=\text{e}^x$}
\end{pspicture*}

\end{center}

\medskip

\begin{minipage}{0.9\linewidth}
\begin{alterqcm}
\AQquestion{Question}{{Proposition 1},{Proposition 2},{Proposition 3}}
\AQquestion{Question}{{Proposition 1},{Proposition 2},{Proposition 3}}
\AQquestion{Question}{{Proposition 1},{Proposition 2},{Proposition 3}}
\AQquestion{Question}{{Proposition 1},{Proposition 2},{Proposition 3}}
\end{alterqcm}
\end{minipage}


\begin{minipage}{0.95\linewidth}
\begin{alterqcm}[VF,lq=125mm,correction,
symb = \dingsquare,
corsymb = \dingchecksquare]
\AQquestion[br=1]{Pour tout $x \in ]-3~;~2],~f'(x) \geqslant 0$.}
\AQquestion[br=2]{La fonction $F$ présente un maximum en $2$}
\AQquestion[br=2]{$\displaystyle\int_{0}^2 f'(x)\:\text{d}x = - 2$}
\end{alterqcm}
\end{minipage}

\end{enumerate}
%%%%%%%%%%%%%%%%%%%%%%%%%%%%%%%%%%%%%%%%%%%%%%%%%%%%%%%%%%%%%%%%%%%%%%%%%
%%%%%%%%%%%%%%%%%%%%%%%%%%%%%%%%%%%%%%%%%%%%%%%%%%%%%%%%%%%%%%%%%%%%%%%%%
\newpage
\subsection{Intervalle de fluctuation ou intervalle de confiance ?}


\begin{center}
\includegraphics[width=0.7\textwidth]{image}
\end{center}

On s'intéresse à une population, dont on étudie un caractère particulier.

\medskip

\begin{bclogo}[couleur = gray!30 , arrondi = 0.1 ,logo = \bclampe , barre = snake , tailleOndu = 1.5]{Deux cas}
\begin{itemize}
\item[$\bullet$] \textbf{\'Echantillonnage}

On utilise un \textbf{intervalle de fluctuation} quand :
	\begin{itemize}
	\item on \textbf{connaît la proportion} $p$ de présence du caractère étudié dans la population,
	\item OU, on \textbf{fait une hypothèse sur la valeur de cette proportion} (on est alors dans le cas de la \og prise de décision\fg~ .
	\end{itemize}

\item[$\bullet$] \textbf{Estimation}

On utilise un \textbf{intervalle de confiance} quand :

on \textbf{ignore la valeur de la proportion} $p$ de présence du caractère dans la population, et on ne formule pas d'hypothèse sur cette valeur.
\end{itemize}
\end{bclogo}

\bigskip
En appliquant l'agorithme de Dijsktra, on obtient :

\begin{center}
\begin{tabular}{|*{8}{>{\centering}m{1.5cm}|}c|}
 \hline
 A&B&C&D&E&F&G&H&Sommet\\
 \hline
 0&$\infty$&$\infty$&$\infty$&$\infty$&$\infty$&$\infty$&$\infty$&A\\
 \hline
  &0+2\par 2(A)&0+1\par 1(A)&0+4\par 4(A)&$\infty$&$\infty$&$\infty$&$\infty$&C\\
 \hline
  &2(A) & &4(A)&1+5\par 6(C)&$\infty$&$\infty$&$\infty$&B\\
 \hline
  & & &4(A)&6(C)&$\infty$&$\infty$&$\infty$&D\\
 \hline
  & & & &6(C)&4+1\par 5(D) &$\infty$&$\infty$&F\\
 \hline
  & & & &6(C)& &5+4\par 9(F)&5+1\par 6(F)&E\\
 \hline
  & & & & & &6+2\par 8(E)&6(F)&H\\
 \hline
  & & & & & &6+4\par 8(E)& &G\\
 \hline
 \end{tabular}
\end{center}



\[\intervalleff{\frac{5}{2}}{-2}\]


\textbf{\textit{Exemples :}}
\begin{enumerate}
\item On dispose d'une pièce de monnaie.

Comment décider qu'elle est \og équilibrée\fg~ ou pas ?

On va ici faire l'hypothèse que la fréquence d'apparition de \og Pile\fg~, par exemple, est égale à 0,5, et on va tester cette hypothèse.

On est dans une \textbf{situation d'échantillonnage}.

\item Une usine fabrique des fusées de feux d'artifice.

Sur 100 fusées choisies au hasard à l'issue du processus de fabrication et mises à feu, on trouve 12 fusées qui ne fonctionnent pas.

Comment se faire une idée de la proportion de fusées défectueuses dans la production ?

On est dans une \textbf{situation d'estimation} : on n'a, au départ, aucune idée de la valeur de la proportion étudiée dans la production.
\end{enumerate}

\medskip

\begin{center}
\PapierMill{10cm}{10cm}
\end{center}

\begin{bclogo}[couleur = blue!30 , arrondi = 0.1 ,logo = \bclampe , barre = snake , tailleOndu = 1.5]{Exercice :}


\textbf{Texte à trous :}

Un enfant a réalisé $\np{4040}$ \trou{lancers} d'une pièce de monnaie, et il a obtenu $\np{2048}$ fois le résultat \og Pile\fg~.

Nous allons nous intéresser au problème suivant : \og Peut-on considérer que la pièce utilisée est équilibrée ?\fg~

\begin{enumerate}
\item Prendre une décision en utilisant l'intervalle de fluctuation au seuil de 95\% sous la forme étudiée en classe de 2\up{nde}.
\item Reprendre le même travail, en utilisant un intervalle de fluctuation au seuil 95\% obtenu à partir de la loi binomiale vue en classe de 1\up{ère}.
\end{enumerate}
\end{bclogo}


\newpage

\textbf{\textit{Solution :}}

\begin{enumerate}
\item En classe de 2\up{nde}, on a vu que pour $n>25$ et $0,2<p<0,8$, l'intervalle $\left[p-\frac{1}{\sqrt{n}};p+\frac{1}{\sqrt{n}}\right]$ constitue un intervalle de fluctuation au seuil 95\%.

Ici, si on appelle $p$ la \og proportion théorique\fg~ de résultats \og Pile\fg~ obtenus sur un très grand nombre de lancers, on veut tester l'hypothèse selon laquelle la pièce est équilibrée, c'est-à-dire l'hypothèse $p=0,5$.

Pour l'échantillon considéré, on a $n=\np{4040}$ ; les conditions d'application sont donc réalisées, et on obtient l'intervalle de fluctuation 
\[\left[0,5-\frac{1}{\sqrt{\np{4040}}};0,5+\frac{1}{\sqrt{\np{4040}}}\right]\]
En arrondissant à $10^{-4}$ près, on obtient l'intervalle de fluctuation $I_1=[\np{0,4843};\np{0,5157}]$.

La fréquence observée dans l'échantillon est donnée par :
\[f=\frac{\np{2048}}{\np{4040}}\approx \np{0,5069}~\text{à}~10^{-4}~\text{près}\]

Ainsi, on a $f\in I_1$, donc on accepte l'hypothèse de pièce équilibrée au seuil de confiance de 95\%.

\item En supposant que la pièce est équilibrée, la variable aléatoire $X$ qui dénombre les résultat \og Pile\fg~ obtenus parmi les $\np{4040}$ lancers réalisés suit une loi binomiale de paramètres $n=\np{4040}$ et $p=0,5$.

On a vu en classe de 1\up{ère} que l'ntervalle de fluctuation à 95\% associé à la variable aléatoire $X$ est l'intervalle $I_2=\left[\frac{k_1}{n};\frac{k_2}{n}\right]$, où $k_1$ est le plus petit entier $k$ vérifiant $P(X\leq k)> 0,025$ et $k_2$ le plus petit entier $k$ vérifiant $P(X\leq k)\geq 0,975$.

On détermine $k_1$ et $k_2$ à l'aide d'un tableur (ou d'une calculatrice...mais le tableau est très long, ici).

On obtient $k_1=\np{1958}$ et $k_2=\np{2082}$, qui donnent comme intervalle de fluctuation :
\[I_2=\left[\frac{\np{1958}}{\np{4040}};\frac{\np{2082}}{\np{4040}}\right]\]
soit en arrondissant à $10^{-4}$ près $I_2=[\np{0,4847};\np{0,5153}]$.

La fréquence observée $f\approx \np{0,5069}$ vérifie $f\in I_2$.

On est ici aussi conduit à accepter l'hypothèse de pièce équilibrée au seuil de confiance de 95\%.
\end{enumerate}
%%%%%%%%%%%%%%%%%%%%%%%%%%%%%%%%%%%%%%%%%%%%%%%%%%%%%%%%%%%%%%%%%%%%%%%%%

\newpage
%%%%%%%%%%%%%%%%%%%%%%%%%%%%%%%%%%%%%%%%%%%%%%%%%%%%%%%%%%%%
\textbf{\textsc{Exercice 2 :}\hfill 8 points}% ex 109 p 29 indice
%%%%%%%%%%%%%%%%%%%%%%%%%%%%%%%%%%%%%%%%%%%%%%%%%%%%%%%%%%%%
\medskip

Une agence de voyages propose un circuit touristique pour visiter trois villes A, B et C.

Le client peut choisir la durée de son séjour dans chacune des villes.

L'agence propose des tarifs qui différent selon la période. Il existe trois périodes touristiques : une période haute, une période moyenne et une période basse.

Les prix journaliers, en \textbf{centaines d'euros} par personne, dans les différents lieux sont donnés dans le tableau suivant :

\begin{center}
\newcolumntype{K}[1]{>{\columncolor{#1}}l}
\begin{tabular}{|K{gray}|*{3}{c|}}
\cline{2-4}
\multicolumn{1}{c|}{} & \cellcolor{yellow} \bf Ville A & \cellcolor{yellow}\bf Ville B & \cellcolor{yellow}\bf Ville C\\
\hline
\bf Période haute & 2,5 & 3,5 & 1,5\\
\hline
\bf Période moyenne & 2 & 2,2 & 1,5\\
\hline
\bf Période basse & 1,5 & 1,2 & 1\\
\hline
\end{tabular}
\end{center}

On considère la matrice $M=\begin{pmatrix}
2,5&3,5&1,5\\2&2,2&1,5\\1,5&1,2&1
\end{pmatrix}$

\begin{enumerate}[\bf 1.]
\item M. Gibbs choisit un circuit de trois jours : deux jours dans la ville A et un jour dans la ville C.

Calculer, à l'aide d'un \textbf{produit de matrices}, le prix de son séjour pour chaque période.

\item La matrice ci-dessous donne toutes les combinaisons d'un circuit de trois jours :

\[D=\begin{pmatrix}
3&2&2&1&1&1&0&0&0&0\\
0&1&0&2&0&...&...&...&...&...\\
0&0&1&0&2&...&...&...&...&...
\end{pmatrix}\]

Ainsi, la troisième colonne indique que M. Gibbs choisit un circuit de deux jours dans la ville A et un jour dans la ville C.

	\begin{enumerate}[\bf a.]
	\item Compléter la matrice $D$.
	
	\item Calculer la matrice $M\times D$.
	
	\end{enumerate}
\item En déduire quels sont les circuits de trois jours possibles pour que le séjour coûte moins de :
	\begin{enumerate}[\bf a.]
	\item 700 \euro~ en période haute ;
	
	\item 560 \euro~ en période moyenne ;
	
	\item 330 \euro~ en période basse.
	
	\end{enumerate}

\item Dans une publicité, l'agence de voyages affirme qu'un circuit complet de 10 jours est possible au prix de \np{2800} \euro~ en période haute, \np{2000} \euro~ en période moyenne et \np{1250} \euro~ en période basse.

Comment se compose ce voyage ?

\end{enumerate}



%%%%%%%%%%%%%%%%%%%%%%%%%%%%%%%%%%%%%%%%%%%%%%%%%%%%%%%%%%%%%%%%%%%%%%%%%
\newpage
\section{Intervalles de fluctuation}

\subsection{Intervalle de fluctuation asymptotique}

On dispose d'un urne contenant un très grand nombre de boules rouges et bleues.

On sait que la proportion de boules rouges dans l'urne est égale à $p=0,4$.

Si on tire, successivement avec remise, $n$ boules dans l'urne ($n$ entier naturel non nul), et si on appelle $X_n$ la variable aléatoire dénombrant les boules rouges tirées, alors $X_n$ suit une loi binomiale $\scr{B}(n;p)$.

\medskip

\begin{bclogo}[couleur = gray!30 , arrondi = 0.1 ,logo = \bclampe , barre = snake , tailleOndu = 1.5]{Théorème et définition}

Soit $X_n$ une variable aléatoire suivant une loi binomiale $\scr{B}(n;p)$, et $\alpha$ un réel tel que $0<\alpha<1$.

Si $X$ est une variable aléatoire suivant la loi normale centrée réduite $\scr{N}(0;1)$, on appelle $u_{\alpha}$ l'unique réel tel que :
\[P(-u_{\alpha}\leq X\leq u_{\alpha})=1-\alpha\]
\begin{itemize}
\item[$\bullet$] On appelle $I_n$ l'intervalle :
\[I_n=\left[p-u_{\alpha}\frac{\sqrt{p(1-p)}}{\sqrt{n}};p+u_{\alpha}\frac{\sqrt{p(1-p}}{\sqrt{n}}\right]\]
Alors,
\[\limn P\left(\frac{X_n}{n}\in I_n=1-\alpha\right)\]

\item[$\bullet$] L'intervalle $I_n$ contient la fréquence $F_n=\frac{X_n}{n}$ avec une probabilité qui se rapproche de $1-\alpha$ lorsque $n$ augmente : 

on dit que c'est l'\textbf{intervalle de fluctuation asymptotique de $F_n$ au seuil $1-\alpha$}.
\end{itemize}
\end{bclogo}

\bigskip

\begin{bclogo}[couleur = yellow!30 , arrondi = 0.1 ,logo = \bclampe , barre = snake , tailleOndu = 1.5]{Démonstration (exigible) :}

On pose $Z_n=\frac{X_n-np}{\sqrt{np(1-p)}}$ et on applique le théorème de Moivre-Laplace.

Si $X$ suit la loi normale $\scr{N}(0;1)$ : $\limn P\left(Z_n\in [-u_{\alpha};u_{\alpha}]\right)=P(X\in [-u_{\alpha};u_{\alpha}]=1-\alpha$.

Or, $Z_n\in [-u_{\alpha};u_{\alpha}] \Lr -u_{\alpha}\leq \frac{X_n-np}{\sqrt{np(1-p)}}\leq u_{\alpha} \Lr np-u_{\alpha}\sqrt{np(1-p}\leq X_n \leq np+u_{\alpha}\sqrt{np(1-p)}$.

$Z_n\in [-u_{\alpha};u_{\alpha}] \Lr p-u_{\alpha}\frac{\sqrt{p(1-p}}{\sqrt{n}}\leq \frac{X_n}{n}\leq p+u_{\alpha}\frac{\sqrt{p(1-p)}}{\sqrt{n}}$.

Donc, $Z_n\in [-u_{\alpha};u_{\alpha}] \Lr \frac{X_n}{n}\in I_n$.
\end{bclogo}

\bigskip

\textbf{\textit{Exemple :}}

On tire 50 boules de l'urne décrite ci-dessus, et on souhaite déterminer un intervalle de fluctuation au seuil 0,9 (c'est-à-dire $\alpha=0,1$).

À l'aide de la calculatrice, on trouve 1,645 comme valeur approchée à $10^{-3}$ près de $u_{0,1}$.

On obtient donc comme intervalle de fluctuation :
\[I_{50}=\left[0,4-u_{0,1}\frac{\sqrt{0,4\times 0,6}}{\sqrt{50}};0,4+u_{0,1}\frac{\sqrt{0,4\times 0,6}}{\sqrt{50}}\right]\]
soit $I_{50}=[0,286;0,514]$.

Ainsi, en effectuant 50 tirages dans cette urne, la fréquence d'apparition d'une boule rouge est comprise entre 0,28 et 0,514 avec une probabilité d'environ 0,9.

Pour 500 tirages, au même seuil de 0,9, on obtient $I_{500}=[0,364;0,436]$.

\medskip

\textbf{\textit{Remarque :}}

La longueur de l'intervalle a, pour un même seuil, été divisée par plus de 3 en passant de 50 à 500 tirages \quad !!

\medskip

\begin{bclogo}[couleur = gray!30 , arrondi = 0.1 ,logo = \bclampe , barre = snake , tailleOndu = 1.5]{Propriété}

L'intervalle de fluctuation asymptotique au seuil de 95\% pour une variable aléatoire $X_n$ suivant une loi binomiale $\scr{B}(n;p)$ est l'intervalle :
\[I_n=\left[p-1,96\frac{\sqrt{p(1-p)}}{\sqrt{n}};p+1,96\frac{\sqrt{p(1-p)}}{\sqrt{n}}\right]\]

\end{bclogo}

\bigskip

\textbf{\textit{Démonstration :}}

Cette propriété découle de celle vue au chapitre 11 concernant la loi normale.

On a vu en effet que $u_{0,05}\approx 1,96$.

\medskip

\textbf{\textit{Remarque :}}

En majorant $1,96\sqrt{p(1-p)}$ par 1, on retrouve l'intervalle de fluctuation présenté en classe de 2\up{nde}.

La majoration précédente résulta de ce qui suit :

$p(1-p)=-p^2+p$ est un trinôme du second degré.

Il admet un maximum en $p=0,5$, qui est égal à 0,25.

Ainsi, pour tout $p\in [0;1], 0\leq p(1-p)\leq 0,25$.

La fonction racine carrée étant croissante sur $\R_+$, on en déduit : $0\leq \sqrt{p(1-p)}\leq 0,5$, 

puis $0\leq 1,96\sqrt{p(1-p)}\leq 1,95\times 0,5\leq 2\times 0,5=1$.

%%%%%%%%%%%%%%%%%%%%%%%%%%%%%%%%%%%%%%%%%%%%%%%%%%%%%%%%%%%%%%%%%%%%%%%%%
%%%%%%%%%%%%%%%%%%%%%%%%%%%%%%%%%%%%%%%%%%%%%%%%%%%%%%%%%%%%%%%%%%%%%%%%%
\subsection{Exercice}

\begin{bclogo}[couleur = blue!30 , arrondi = 0.1 ,logo = \bclampe , barre = snake , tailleOndu = 1.5]{\'Enoncé}

Dans un casino, il a été décidé que les \og machines à sous\fg~ doivent être réglées sur une fréquence de gain du joueur de $g=0,06$.

Une fréquence inférieure est supposée faire \og fuir \fg~ le client, et une fréquence supérieure est susceptible de ruiner le casino.

Trois contrôleurs différents vérifient une même machine.

Le premier a joué 50 fois et gagné 2 fois,

le second a joué 120 fois et gagné 14 fois,

le troisième a joué 400 fois et gagné 30 fois.

En utilisant des intervalles de fluctuation asymptotiques au seuil 95\%, examiner dans chaque cas la décision à prendre par le contrôleur, à savoir accepter ou rejeter l'hypothèse $g=0,06$.
\end{bclogo}

\bigskip

\textbf{\textit{Solution :}}

\begin{itemize}
\item[$\bullet$] \textbf{Premier contrôleur}
Le contrôle a porté sur 50 parties, donc $n=50$.

On a bien $n\geq 30$, mais $np=50\times 0,06=3$ et donc $np<5$.

On n'est pas dans les conditions d'application d'un intervalle de fluctuation asymptotique.

Ce contrôle ne peut rien donner de probant en termes de prise de décision.

\item[$\bullet$] \textbf{Deuxième contrôleur}

Le contrôle a porté sur 120 parties, donc $n=120, np=7,2$ et $n(1-p)=112,8$.

Les conditions d'utilisation d'un intervalle de fluctuation asymptotique sont donc réunies.

L'intervalle de fluctuation asymptotique au seuil 95\% pour un échantillon de taille $n=120$ est égal à :
\[I_{120}=\left[0,06-1,96\frac{\sqrt{0,06\times 0,94}}{\sqrt{120}};0,06+1,96\frac{\sqrt{0,06\times 0,94}}{\sqrt{120}}\right]\]
soit en arrondissant les bornes à $10^{-4}$ près : $I_{120}=[\np{0,0175};\np{0,1025}]$.

La fréquence observée par le deuxième contrôleur est $f_2=\frac{14}{120}\approx \np{0,1167}$.

Cette fréquence est à l'extérieur de l'intervalle de fluctuation $I_{120}$, et le second contrôleur est conduit à rejeter l'hypothèse $g=0,06$.

\item[$\bullet$] \textbf{Troisième contrôleur}

La contrôle a porté sur 400 parties, donc $n=400, np=24$ et $n(1-p)=376$.

Les conditions d'utilisation d'un intervalle de fluctuation asymptotique sont donc réunies.

L'intervalle de fluctuation asymptotique au seuil 95\% pour un échantillon de taille $n=400$ est égal à :

\[I_{400}=\left[0,06-1,96\frac{\sqrt{0,06\times 0,94}}{\sqrt{400}};0,06+1,96\frac{\sqrt{0,06\times 0,94}}{\sqrt{400}}\right]\]
soit en arrondissant les bornes à $10^{-4}$ près : $I_{400}=[\np{0,0367};\np{0,0833}]$.

La fréquence observée par le deuxième contrôleur est $f_3=\frac{30}{400}= \np{0,075}\in I_{400}$.

Ce dernier est donc amené à accepter l'hypothèse $g=0,06$.

\end{itemize}

\medskip

\begin{bclogo}[couleur = gray!30 , arrondi = 0.1 ,logo = \bclampe , barre = snake , tailleOndu = 1.5]{Remarques}

\begin{enumerate}
\item Pour tester une hypothèse à l'aide d'un intervalle de fluctuation asymptotique, on doit vérifier ses conditions d'utilisation :
\[n\geq 30, \quad np\geq 5, \quad n(1-p)\geq 5\]

Ces conditions vérifiées, on détermine l'intervalle de fluctuation $I_n$ au seuil demandé (souvent 95\%).

Si la fréquence observée $f$ n'appartient pas à $I_n$, on rejette l'hypothèse.

\item L'intervalle de fluctuation étudié est \og bilatéral\fg~ : on est conduit à rejeter l'hypothèse dans le cas où la fréquence observée est \og trop grande\fg~, ou \og trop petite \fg~.
\end{enumerate}
\end{bclogo}

%%%%%%%%%%%%%%%%%%%%%%%%%%%%%%%%%%%%%%%%%%%%%%%%%%%%%%%%%%%%%%%%%%%%%%%%%


%%%%%%%%%%%%%%%%%%%%%%%%%%%%%%%%%%%%%%%%%%%%%%%%%%%%%%%%%%%%

\end{spacing}

%%%%%%%%%%%%%%%%%%%%%%%%%%%%%%%%%%%%%%%%%%%%%%%%%%%%%%%%%%%%
%%%%%%%%%%%%%%%%%%%%%
%% FIN DU DOCUMENT %%
%%%%%%%%%%%%%%%%%%%%%
\end{document}